\section{Implementation of the smearing}
\label{sec:implement}

To implement this analytic link smearing scheme, the efficient evaluation
of $\exp( i Q)$, where $Q$ is a traceless Hermitian $3\times 3$ matrix,
is required. The Cayley-Hamilton theorem states that every matrix is a zero
of its characteristic polynomial, so that
\begin{equation}
  Q^3 - c_1\ Q - c_0\ I = 0,
\label{eq:CayHam}
\end{equation}
where
\begin{eqnarray}
 c_0 &=&  \det Q=\textstyle\frac{1}{3}\Tr(Q^3),\\
 c_1 &=& \textstyle\frac{1}{2} \Tr(Q^2)\geq 0.
\end{eqnarray}
The Hermiticity of $Q$ requires $27c_0^2-4c_1^3\leq 0$  and
the definition of $Q$ given in Eq.~(\ref{eq:Qdef}) restricts the possible
values of $c_1$.  Thus, the coefficients $c_0$ and $c_1$ satisfy
\begin{equation}
 -c_0^{\rm max}\leq c_0 \leq c_0^{\rm max},\qquad
0 \leq c_1 \leq c_1^{\rm max},
\end{equation}
where
\begin{eqnarray}
 c_0^{\rm max}&=&2\left(\frac{c_1}{3}\right)^{3/2},\\
  c_1^{\rm max} &=& \frac{1}{32}
  (69+11\sqrt{33})\Bigl(\displaystyle\sum_{\nu\neq\mu}
\rho_{\mu\nu}\Bigr)^2,
\end{eqnarray}
for each $\mu$.  Eq.~(\ref{eq:CayHam}) implies that $Q^n$ for integer
$n\geq 3$ can be expressed in terms of $Q^2$, $Q$, and the identity
matrix $I$.  Hence, we can write
\begin{equation}
 e^{i Q} = f_0\ I + f_1\ Q + f_2\ Q^2,
\label{eq:fexpand}
\end{equation}
where the three scalar coefficients $f_j=f_j(c_0,c_1)$ are
basis independent, depending only on $c_0$ and $c_1$.  Eq.~(\ref{eq:fexpand})
is valid for {\em any} $3\times 3$ Hermitian, traceless matrix $Q$.

Let $q_1,q_2,q_3$ denote the three eigenvalues of $Q$. Since $Q$ is
Hermitian and traceless, we know that these are real numbers satisfying
$q_1+q_2+q_3=0$.  These eigenvalues are the three roots of a cubic
polynomial which can be easily determined:
\begin{eqnarray}
q_1 &=& 2u,\\
q_2 &=& -u+w,\\
q_3 &=& -q_1-q_2=-u-w,
\end{eqnarray}
where
\begin{eqnarray}
u &=& \sqrt{\textstyle\frac{1}{3}c_1}
  \ \cos\left(\textstyle\frac{1}{3}\theta\right),\label{eq:udef}\\
w &=& \sqrt{c_1}\ \sin\left(\textstyle\frac{1}{3}\theta\right),\label{eq:wdef}\\
\theta &=& \arccos\left(\frac{c_0}{c_0^{\rm max}}\right).\label{eq:thetadef}
\end{eqnarray}
Given these eigenvalues, the matrix $Q$ can be written
\begin{equation}
 Q = M\ \Lambda_Q\ M^{-1},\qquad 
\Lambda_Q= \left[\begin{array}{ccc} q_1 &0 & 0\\ 0 & q_2 & 0\\
 0 & 0 & q_3\end{array}\right] ,
\end{equation}
where $M$ is a unitary matrix.  Then it easily follows that
\begin{equation}
  e^{i \Lambda_Q}=f_0 I + f_1\Lambda_Q + f_2\Lambda_Q^2.
\end{equation}
Explicitly, we have the following linear system of equations to solve:
\begin{equation}
 \left[\begin{array}{ccc} 1 & q_1 & q_1^2 \\
  1 & q_2 & q_2^2 \\  1 & q_3 & q_3^2 
 \end{array}\right]\left[\begin{array}{c}
  f_0\\f_1\\f_2\end{array}\right]=\left[\begin{array}{c}
 e^{i q_1}\\e^{i q_2}\\e^{i q_3}\end{array}\right].
\label{eq:system}\end{equation}
If all three eigenvalues are distinct, this system of equations
has a unique solution, but when two of the eigenvalues are exactly the
same, the solution to Eq.~(\ref{eq:system}) is not unique and one of the
$f_j$'s can be freely set to any value. The case of degenerate eigenvalues
occurs when $c_0\rightarrow \pm c_0^{\rm max}$.  

In practice, it is extremely unlikely that an exact degeneracy will be
encountered during a numerical simulation.  Much more likely is the
possibility that two eigenvalues are nearly, but not quite, equal.
Encounters with both near and exact degeneracies can be handled
by expressing the $f_j$ in terms of $u$
and $w$ in order to isolate and tame the numerically sensitive part.
One finds that the $f_j$ factors can be written
\begin{equation}
   f_j = \frac{h_j}{(9u^2-w^2)},
\label{eq:fdef1}
\end{equation}
where the $h_j$ are well-behaved functions given by
\begin{eqnarray}                                           
  h_0 &=& (u^2\!-\!w^2)e^{2i u}+e^{-i u}\bigl\{ 8u^2\cos( w)\nonumber\\
  &&\quad + 2iu(3u^2\!+\!w^2)\xi_0( w)\bigr\}, \\
  h_1 &=& 2ue^{2i u}-e^{-i u}\bigl\{ 2u\cos( w) \nonumber\\
  &&\quad - i(3u^2\!-\!w^2)\xi_0( w)\bigr\}, \\
  h_2 &=& e^{2i u}-e^{-i u}\left\{ \cos( w)
  + 3iu\xi_0( w)\right\},
\label{eq:hdef}
\end{eqnarray}
defining
\begin{equation}
 \xi_0(w) = \frac{\sin w}{w}.
\end{equation}
The $w\rightarrow 0$ problem as $c_0\rightarrow c_0^{\rm max}$
has been completely contained in $\xi_0$.  To numerically evaluate 
this factor, one uses, for example,
\[
 \xi_0(w) = \left\{ \begin{array}{ll}
  1\!-\!\frac{1}{6}w^2\left(1\!-\!\frac{1}{20}w^2
  \left(1\!-\!\frac{1}{42}w^2\right)\right) ,
 & \  \vert w\vert\leq 0.05,\\[2mm]
  \sin(w)/w, & \ \vert w\vert>0.05.
 \end{array}\right.
\]
However, the $w\rightarrow 3u\rightarrow \sqrt{3}/2$ limit as
$c_0\rightarrow -c_0^{\rm max}$ has not been tamed.
Fortunately, this problem can be circumvented using the following
symmetry relations under $c_0\rightarrow -c_0$:
\begin{equation}
   f_j(-c_0,c_1)=(-)^j\ f_j^\ast(c_0,c_1). \label{eq:fsym}
\end{equation}
Thus, the determination of the $f_j$ coefficients for $c_0<0$ can
be achieved by computing the coefficients for $\vert c_0\vert$
and utilizing Eq.~(\ref{eq:fsym}).  This means that a
situation in which the denominator in Eq.~(\ref{eq:fdef1})
becomes nearly zero will never be encountered.  It should be
emphasized that the $f_j$ coefficients are smooth, non-singular
functions of $c_0$ and $c_1$.  There are no actual singularities
as $c_0\rightarrow\pm c_0^{\rm max}$.  The numerical evaluation
of the $f_j$ coefficients to machine precision in these limits
simply requires additional care.

As an aside, if one writes $Q=\frac{1}{2}\sum_{a=1}^8 Q_a\lambda^a$
in terms of the eight Gell-Mann matrices $\lambda^a$, then
\begin{eqnarray}
  e^{i Q}&=&u_0+\textstyle\frac{1}{2}\sum_{a=1}^8 u_a \lambda^a,\\
  u_0 &=& f_0+\textstyle\frac{2}{3}c_1f_2,\\
  u_a &=& f_1 Q_a + \textstyle\frac{1}{2} f_2 d^{abc}Q_bQ_c,\\
  c_0 &=& \textstyle\frac{1}{12}d^{abc}Q_a Q_bQ_c,\\
  c_1 &=& \textstyle\frac{1}{4}Q_aQ_a,
\end{eqnarray}
where $d^{abc}$ and $f^{abc}$ are the real symmetric and antisymmetric
structure constants, respectively.


